\chapter{Methodology}
\label{chapter:methodology}


The methodology adopted in this thesis treats decoding as a two-stage, quality-first process that is shared by both automatic speech recognition and image captioning. For each input, the system first generates a set of candidate sequences and then selects a single output according to utilities that better reflect human-aligned quality than raw likelihood. Throughout the experiments, the number of candidates and the decoding controls that modulate diversity are varied, and both quality and computational cost are recorded in order to study how performance scales with the size and composition of the candidate set.

\section{Tasks and Data}

The speech experiments use LibriSpeech as the primary English benchmark, focusing on splits where per-utterance audio files and reference transcripts are available. To probe robustness in European Portuguese, complementary experiments can be conducted on CAM\~{O}ES, with decoding configured for the target language. Candidate transcripts are produced by compact and mid-sized recognizers so that scaling in the number of hypotheses remains practical in wall-clock time; in particular, Whisper-family encoder--decoder models provide a convenient generator for controlled candidate generation and sampling-based variation, and NeMo Conformer/FastConformer ASR models can be used as an additional point of comparison under beam-style decoding.

The captioning experiments operate on MS~COCO, where five human references per image support stable multi-reference evaluation. A strong pretrained captioning model is used as the base generator, and candidate diversity is induced through the same family of decoding controls considered in speech, with the difference that reference-free utilities can additionally exploit image grounding.

\section{Candidate Generation}

Candidate generation follows the same general pattern across tasks. A likelihood-oriented baseline is obtained through beam search with length normalization, returning either a single best hypothesis or an \(N\)-best list. To obtain less correlated candidate sets, diversity-promoting variants are considered. Diverse beam partitions the beam into groups and applies a penalty that discourages later groups from reusing the same next-token choices made by earlier groups, thereby exploring alternative continuations. Stochastic decoding is induced through sampling controls, namely nucleus (top-\(p\)) sampling and temperature, which reshape the next-token distribution before sampling so that the candidate pool covers multiple plausible regions of the model distribution. For each input, candidates from one or more strategies can be pooled and then deduplicated under a simple text normalization rule, so that the evidence set is not dominated by near-identical hypotheses.

\section{Utilities and Evaluation Measures}

Evaluation relies on both reference-based and reference-free metrics, and the same measures can serve as utilities in the second-stage selection. In speech, the primary objective is word error rate. If \(S\), \(D\), and \(I\) denote substitutions, deletions, and insertions needed to transform a hypothesis transcript \(h\) into a reference transcript \(r\), and \(N\) is the number of reference words, then
\begin{equation}
\mathrm{WER}(h,r) \;=\; \frac{S + D + I}{N}.
\end{equation}
When reporting results as percentages, the value is multiplied by \(100\%\). Character error rate (CER) is used as a complementary edit-distance measure when appropriate. To capture meaning preservation beyond literal edits, SeMaScore is reported alongside WER as a semantic-oriented view of transcript quality. In scenarios where references are unavailable or when selection must be performed without consuming references, a reference-free proxy such as NoRefER can be used as a monitoring signal and, when desired, as a quality estimator for ranking competing hypotheses.

In captioning, the standard multi-reference COCO protocol emphasizes content and semantic adequacy, and reference-based metrics can be used for reporting and, when appropriate, as utilities for selection. When references are sparse or absent, reference-free measures become central: CLIPScore quantifies image--text alignment directly and enables analysis and selection grounded in the input image rather than in overlap with textual references.

\section{Selection Rules and Risk-Based Reranking}

The second stage selects a single output from the candidate pool by applying decision rules that target quality rather than likelihood alone. A direct baseline is \(n\)-best selection under a model score, which approximates a likelihood-oriented choice within a restricted candidate set. Quality-aware selection replaces that objective with utilities designed to reflect the evaluation target. In its simplest form, the system selects the candidate that maximizes a quality estimate applied independently to each hypothesis. In its risk-based form, selection follows a minimum-risk criterion that aggregates evidence across candidates.

Given a candidate pool \(\mathcal{Y}=\{y_1,\dots,y_N\}\) and a utility \(u\), the minimum-risk decision is instantiated over an evidence set \(\mathcal{E}\), typically taken to be the pool itself (or a truncated subset when efficiency requires it), by choosing the hypothesis with the highest average utility against the evidence:
\[
\hat{y} \;=\; \arg\max_{y\in \mathcal{Y}}\; \frac{1}{|\mathcal{E}|}\sum_{y'\in\mathcal{E}} u(y,y').
\]
In speech, this criterion is naturally instantiated with edit-distance utilities by taking \(u\) to be the negative of an error measure (e.g., \(-\mathrm{WER}\) or \(-\mathrm{CER}\)) computed between pairs of hypotheses, which encourages consensus-like transcripts that are close, on average, to other plausible outputs. When semantic behavior is of interest, utilities based on SeMaScore can be used in a reference-based setting to favor meaning-preserving hypotheses beyond literal word matches. In quality-aware variants that operate without references, the evidence aggregation can be paired with weights derived from a reference-free quality estimator, so that the risk emphasizes agreement with hypotheses that are predicted to be higher quality.

In captioning, when multi-reference data are available, utilities derived from COCO metrics can be used to favor captions that are strong under the target evaluation criterion; when references are not used, selection can be grounded by image--text alignment through a CLIP-based score, either as a direct per-candidate utility or as a mechanism to bias risk-based selection toward visually grounded candidates. Across both domains, the essential methodological point is that selection is defined over complete sequences under utilities intended to capture the desired notion of quality, rather than by maximizing tokenwise likelihood alone.

\section{Scaling Protocol and Computational Measurement}

The experimental protocol is designed to expose scaling trends and practical trade-offs. For each task, model, and decoding configuration, the candidate count \(N\) is swept and corpus-level quality is compared against compute. Compute is tracked by measuring end-to-end decoding time per input under consistent hardware settings, and by reporting how runtime changes as additional candidates are generated and evaluated. As \(N\) increases, improvements are expected to be steeper when candidate pools are less correlated, since additional hypotheses contribute new evidence; conversely, gains saturate as added hypotheses become redundant. This protocol makes it possible to quantify regimes where computation is better spent on generating and selecting among more diverse candidates, rather than on relying solely on a single likelihood-first decode, and to document diminishing returns beyond moderate candidate set sizes.
